\c"He appeared to them alive again the third day, as the divine
prophets had foretold."--\textit{Josephus}.


How often have men missed the sight of great historic occurrences,
in their attention to the routine of life!  So it was that Quintus
did not witness the tragic events of that Passover week on which
human destiny was to turn.  To Tyre on the Great Sea he had gone,
to arrange for the landing of a new quota of troops from
Brundisium.  The commander at Scopus had chosen him for the
responsible mission, in token of his especial fitness.  The
compliment was pleasing.  But in his absence he was ever thinking
of the promise made by the Teacher in Solomon's Porch, that the
sheep who followed him should have eternal life.

Astir was all Jerusalem, when the knight returned to Scopus.  It
was on the morning after the Lord's resurrection.  That Roman
centurion who had been at Calvary reviewed for Quintus the fateful
happenings.  With pomp reminding of a Roman triumph the Christ had
entered David's city; after four days Iscariot had betrayed him
with a kiss; for blasphemy Pilatus, the procurator, had sentenced
him to the cross; they had put on him a scarlet robe in mockery;
they had hung him between two robbers on the hill of Golgotha; a
brutal soldier now at Scopus had won by lot his seamless robe, and
was jauntily displaying it as a trophy; an uncanny darkness had
covered the Judaean sky; the soldier Longinus had pierced the
sufferer's side; they had buried the dead Christ in the garden tomb
of the Arimathaean Joseph.  Monumental events were these--all new
to Quintus, but destined to be written indelibly in the calendars
of Christendom.

"More than this," continues the centurion, "an amazing rumor is now
abroad in the city that yesterday the dead Christus awoke from his
sleep and has been five times seen by his amazed disciples.  When I
beheld him yield up the ghost, I hailed his death as that of a
devout man, but little did I think that he was a God and would
return from the tomb.  The report says he has now come back.  On
swift wing the rumor has flown through Jerusalem and even into
Pilate's palace."


Down from the heights of Scopus the hurrying feet of Quintus carry
him to Jerusalem.  Doubts and wonderings and half-beliefs fill his
mind.  What if by any shadow of possibility the prediction of the
strange Teacher has been fulfilled, that he should return from the
dead on the third day?  Finding his way to Joseph's garden, Quintus
stands by an empty sepulcher.  There is a group of wondering
visitors near, and among them is one whose inviting face leads
Quintus to accost him.  Not frightened by the sword and armor of
the Roman knight, but assured by his candid look, the other answers
in the Aramaic which both can speak:

"Johannes is my name.  Till three years ago I was a fisherman, up
on the waters of Gennesaret.  Since then I have been a disciple of
this Man from Galilee.  In his company I have heard surprising
words and have felt a heavenly influence.  He was no ordinary
Teacher.  He was indeed from above."

"Is it true," asks Quintus in breathless words, "that your Master
has risen from the grave?  I have been away in Tyrus.  Now in the
Roman camp on Scopus I have heard that he has come forth from the
sepulcher.  What means such a marvelous report?"

"Yes, it is all true," John answers with his face aglow; "this is
the very sepulcher where our Lord was laid.  Your own sentries kept
guard before the tomb securely sealed.  But on the morning of
yesterday there was a shaking of the earth; some angelic visitants
rolled away the stone door of the grave; and our immortal Christus
came forth again.

"Astounding," Quintus interrupts in a whirl of words; "but did he
make any promise of another life for men, before he was put to
death?"

"He truly did," replies the disciple; "when we had eaten the
Passover supper with him, he spoke a word more marvelous than any
of your Roman teachers has ever uttered.  Into the spirit world he
said he was departing, to make ready a room in the Father's ample
house for those who were his own; and on his return he would take
them to be with himself.  Ever since our sad-hearted band have been
comforting themselves with this last promise in the upper room."

"None of our Roman gods has ever promised such a future." responds
Quintus; "but is this all?"

"No," answers the disciple; "on his cross our Christus spoke again
about another experience for men.  By his side was Dysmas, the
crucified robber, grieving for his faults and asking comfort.  When
the cross pain and thirst were over, our Lord replied, the outlaw
should walk with him among the bowers of the beautiful Paradise
beyond this world's horizon.  It was enough.  In this consolation
the tortured Dysmas passed on, with a smile of peace upon his face."

"Have you more wonders to tell?" presses Quintus, in his eagerness,
while the story of the cross begins to compel his judgment and call
for his heart's surrender.

Then, the consummation!  In ecstatic words John tells of the one
final and overmastering proof, in the thought of the eleven
disciples;

"Greatest of all, we have ourselves seen our Friend again.  Five
times already has he showed himself.  First, Mary of Magdala saw
him under the trees of the garden, and spoke with him; then the
other women met him and fell at his feet; next our fellow disciple
Petros saw him; then two of our band walked with him to outlying
Emmaus, and knew him as he broke bread at the journey's end; and
then last evening, he came to ten of us in the Passover room and
spoke his peace on us.

"Perhaps you have all seen a spectral form which has no real
existence," remonstrates Quintus, while all the time he is yielding
himself to the compelling story.

"It cannot be," responds the convincing John; "there have been too
many witnesses for that.  We have seen the very wound made by the
spear of Longinus; we have heard his familiar voice; we have
received his blessing.  Our number is our evidence; it cannot be
possible that all of us have been deceived.  It is surely he, O
Roman soldier, unless the senses of the women and of ten honest men
are far astray.  No other teacher of the East has ever come back
from the sepulcher.  Look and see for yourself.  Yonder is Joseph's
empty tomb.  The Christus is himself the evidence."


What can Quintus do, in the face of such proof as this?  He returns
to Scopus in wildest tumult.  Little does he say to Aulus, his
chosen friend.  The company of Longinus or the centurion he does
not seek.  The time has come--as it comes to all--when he must
commune with himself, and make the decision confronting every soul
that has heard the resurrection story.